% ---------------------------------------------------------------------------
%  Resumen (200-250 palabras), Abstract y Descriptores.
% ---------------------------------------------------------------------------
\cleardoublepage
\chapter*{Resumen}
\addcontentsline{toc}{chapter}{Resumen}

ComprendiA es una aplicación web educativa que transforma vídeos de YouTube en material
de estudio consultable. A partir de la URL de una grabación, el sistema obtiene su
transcripción —reutilizando los subtítulos de YouTube cuando existen y recurriendo al
modelo Whisper en caso contrario—, la divide en fragmentos y genera un \textit{embedding}
por fragmento que permite la búsqueda semántica sobre el contenido. Sobre esa base, un
proceso de análisis automático produce capítulos navegables y conceptos clave con sus
marcas de tiempo, y un asistente conversacional responde preguntas del alumno mediante una
arquitectura de generación aumentada por recuperación (RAG), manteniendo una memoria corta
del diálogo para resolver referencias implícitas.

La aplicación organiza las clases procesadas en asignaturas. Para reducir el trabajo manual,
incorpora una clasificación automática que sugiere a qué asignatura pertenece cada vídeo
nuevo, primero por coincidencia de canal de YouTube y, si no la hay, por similitud semántica
con las asignaturas existentes, creando una nueva cuando ninguna resulta adecuada.

El sistema se ha construido con un \textit{backend} en Java 21 sobre Quarkus, una base de
datos PostgreSQL con la extensión \texttt{pgvector} y un \textit{frontend} en Angular. Esta
memoria describe la motivación, el estado del arte de las tecnologías implicadas, la
especificación de requisitos, el diseño y la implementación del sistema, las pruebas
realizadas y las líneas de trabajo futuras.

\vspace{1cm}
\section*{Abstract}

% TODO: revisar la traducción con un hablante nativo antes de la entrega.
ComprendiA is an educational web application that turns YouTube videos into searchable study
material. Given a video URL, the system obtains its transcript —reusing YouTube subtitles when
available and falling back to the Whisper model otherwise—, splits it into fragments and
generates one embedding per fragment to enable semantic search over the content. On top of
this, an automatic analysis produces navigable chapters and key concepts with their
timestamps, and a conversational assistant answers the student's questions using a
retrieval-augmented generation (RAG) architecture, keeping a short-term memory of the dialogue
to resolve implicit references.

The application organises processed lectures into subjects. To reduce manual effort, it
includes an automatic classifier that suggests which subject each new video belongs to, first
by matching the YouTube channel and, failing that, by semantic similarity with existing
subjects, creating a new one when none fits.

The system is built with a Java 21 backend on Quarkus, a PostgreSQL database with the
\texttt{pgvector} extension and an Angular frontend. This document covers the motivation,
the state of the art, the requirements specification, the design and implementation, the
testing performed and future work.

\vspace{1cm}
\section*{Descriptores / Keywords}

\noindent\textbf{Descriptores:} transcripción automática, búsqueda semántica, \textit{embeddings},
generación aumentada por recuperación (RAG), tecnología educativa.

\noindent\textbf{Keywords:} automatic transcription, semantic search, embeddings,
retrieval-augmented generation (RAG), educational technology.
